\subsection{Proof for the ``soundness" clause for the answer reduction transformation}  \label{sec:ARsoundness}
As mentioned in~\Cref{sec:Answerreductionproof}, this subsection follows a similar structure as~\cite[Section 10.7]{jiMIPRE2022a}

Fix $\alpha, n \in \bN$, let $\cG_n = (\cX_n, \cA_n,\mu_n, D_n)$ be the $n$th game of $\verifiersequence$, and let $\cG_n^{\Answerred} = (\cX_n^{\Answerred}, \cA_n^{\Answerred},\mu_n^{\Answerred}, D_n^{\Answerred})$ and $\cG_n^{\text{Ora}}  = (\cX_n^{\text{Ora}}, \cA_n^{\text{Ora}},\mu_n^{\text{Ora}}, D_n^{\text{Ora}})$ denote the answer reduction transformation given in~\Cref{sec:Answerreductionproof} and~\Cref{sec:Oracularization}. Given a question label $x \in \cX_n^{\Answerred}$, we write $x = (x^{\text{Ora}}, x^{\text{LDL}}, (x^{\text{game}},x^{\text{LDC}}))$ where each elements corresponds to the ``oracularizable question label", ``SLDT question label", question content for the original game $\cG_n$, and  question content for SLDT from~\Cref{fig:Answerredsample} respectively. When we fix certain question labels, we might omit that portion of the question label for simplicity notation. Furthermore, let $(\mathtt{PCPParameter}_{\alpha}, \mathtt{ComputePCP}_{\alpha})$ be the two Turing machine guaranteed by~\Cref{thm:classicalPCPP}, and let $(m^{\text{ans}},m, g, p) = \mathtt{PCPParameter}_{\alpha}(n)$. 

Before we start giving a proof of the ``soundness clause", we first give a first overview on how the proof goes. To show the ``soundness" clause given in~\Cref{eq:ARsoundnessclause}, it is equivalent to show that there exist a polynomial function $\textbf{t}_{\alpha}^{\Answerred}$ with $\textbf{t}_{\alpha}^{\Answerred} = O(\polylog(n), \poly(\eps))$ such that for model $t \in \{*, co\}$
\begin{equation*}
	\omega^{t}(\cG^{\Answerred}_n) > 1 - \eps \Longrightarrow \omega^t(\cG) > 1 -\mathbf{t}^{\Answerred}_{\alpha}(\eps, n). 
\end{equation*}
Hence, assume that $\omega^{t}(\cG^{\Answerred}_n) > 1  - \eps$, and fix strategy in model $t$ such that $\strategy$ succeed at $\cG^{\Answerred}_n$ with probability at most $\eps$. We show the soundness clause by proving the following:
\begin{enumerate}
	\item By using~\Cref{thm:soundnessQLDT}, we first show that there exist a strategy $\strategy^{\text{poly}}$ for $\cG^{\Answerred}_n$ which consist of the provers first performing hidden measurements and sampled $6+m$ low-individual degree polynomials, and then using these polynomials to pass all the low-individual degree polynomial/simultaneous low-individual degree polynomial test for $\cG_n^{\Answerred}$.
	\item Then we use~\Cref{thm:classicalsoundnessPCPP} on the polynomial generated by $\strategy^{\text{poly}}$ to construct a strategy $\strategy^{\text{Ora}}$ for $\cG^{\text{Ora}}_n$ which succeed with probability at least $1- O(\polylog(n), \poly(\eps))$. 
	\item Finally, we conclude the proof by applying~\Cref{lem:oratransformation} to show that there exist a strategy for $\cG_n$  which succeed with probability at least $1- O(\polylog(n), \poly(\eps))$, thus showing the lemma. 
\end{enumerate}
For simplicity of notations, we work with synchronous strategies in order to show point 1 and 2. For a question pair $(x,y) = \left( (x^{\text{Ora}}, x^{\text{LDL}}, (x^{\text{game}},x^{\text{LDC}} )), (y^{\text{Ora}}, y^{\text{LDL}}, (y^{\text{game}},y^{\text{LDC}}))\right)$ for $\cG^{\Answerred}_n$, we observe that the synchronous question pair for $\cG^{\Answerred}_n$ corresponds to the case where $x^{\text{Ora}} = y^{\text{Ora}}$ and $x^{\text{LDL}} = y^{\text{LDL}}$ which occur with constant probability. This implies that the game $\cG^{\Answerred}_n$ is $\frac{1}{c_b}$-balanced for some constant $\frac{1}{c_b}$, and hence any strategy $\strategy$ for $\cG^{\Answerred}_n$ which succeed with probability $1- \eps$ must be $c_b \cdot  \eps$-synchronous. By~\Cref{thm:Rounding}, we have  
\begin{equation}
	\omega^{t}(\cG^{\Answerred}_n) > 1 - \eps  \Longrightarrow	\omega^{t}_s(\cG^{\Answerred}_n) > 1 - \eps - \textbf{s}^{\text{Rounding}}(c_b\eps) = 1 - \delta_1.
\end{equation}
Hence, fix a synchronous strategy $\strategy = (\cL^2(\alicealg, \tau), \ket{\tau}, \{A_a^x \})$ such that $\omega(\cG^{\Answerred}_n, \strategy) > 1 - c_1 \cdot \delta_1$, where $c_1$ is a sufficiently small constant choose later down the proof.


Define the function $\text{eval}_s^n: \Idpoly(p,m^{\text{ans}},p)^{\times n} \rightarrow \bF_{2^p}^{\times n}$ as $\text{eval}_s^n(\textbf{g}_0, \cdots, \textbf{g}_{n-1}) = (\textbf{g}_0(s),\cdots ,\textbf{g}_{n-1}(s))$. We wish to first show the following claim:
\newcommand{\game}{\text{game}}
\begin{claim} \label{claim:ARpolyPVM}
	For all $(x^{\game},y^{\game}) \in \cX^2$, there exist six sets of PVM in $\alicealg'$, 
	\begin{itemize}
		\item $\{G_{\overline{\textbf{g}}}^{(\text{Prover, A}), x^{\game}}\}_{\overline{\textbf{g}} \in \Idpoly(p,m^{\text{ans}},p)}$,
		\item $\{G_{\overline{\textbf{g}}}^{(\text{Prover, B}), y^{\game}}\}_{\overline{\textbf{g}} \in \Idpoly(p,m^{\text{ans}},p)}$,
		\item $\{G_{\overline{\textbf{g}}}^{(\text{Ora})_{o}, (x^{\game}, y^{\game})}\}_{\overline{\textbf{g}} \in \Idpoly(p,m^{\text{ans}},p)}$ for $o \in \{0,1,2\}$,
		\item $\{G_{\textbf{g}_{U_0}, \cdots, \textbf{g}_{U_4}, \textbf{g}_{\Gamma}, \textbf{g}_{B_0} \cdots,\textbf{g}_{B_{m-1}}}^{(\text{Ora}), (x^{\game}, y^{\game})}\}_{\textbf{g}_v \in \Idpoly(p,m,p)}$, 
	\end{itemize}
	such that the following hold: For $s \in \bF_{2^p}^{m^{\text{ans}}}$ and $n \in \bN$,
	\begin{align}
		\label{eq:ARsoundAlicepoly} A^{(\text{Prover, A}), (\text{Point}), x^{\game} ,s}_{u} &\simeq_{O(\poly(\eps, \log(n), \alpha))} G_{[\textbf{eval}_{s}^1 | u]}^{(\text{Prover, A}), x^{\game}}, \\
		\label{eq:ARsoundBobpoly} A^{(\text{Prover, B}), (\text{Point}), y^{\game} ,s}_{u} &\simeq_{O(\poly(\eps, \log(n), \alpha))} G_{[\textbf{eval}_{s}^1 | u]}^{(\text{Prover, B}), y^{\game}}, \\
		\label{eq:ARsoundvarpoly}  A^{(\text{Ora})_{o}, (\text{Point}), ((x^{\game}, y^{\game}) ,s)}_{u} &\simeq_{O(\poly(\eps, \log(n), \alpha))} G_{[\textbf{eval}_{s}^1 | u]}^{(\text{Ora})_{o}, (x^{\game}, y^{\game})}, \quad o \in \{0,1,2\}, \\
		\label{eq:ARsoundAllpoly} A^{(\text{Ora}), (\text{Point}), ((x^{\game}, y^{\game}) ,s)}_{u_0, \cdots, u_4, \gamma, \beta_0, \cdots,\beta_{m-1}} &\simeq_{O(\poly(\eps, \log(n), \alpha))} G_{[\textbf{eval}_s^{6+m}|( u_0, \cdots, u_4, \gamma, \beta_0, \cdots,\beta_{m-1})]}^{\text{Ora}, (x^{\game}, y^{\game})}, 
	\end{align}
	where $\simeq$ for the above four equations is defined over the distribution $(x^{\game}, y^{\game}) \sim \mu_n$ and $s \sim \bF_{2^p}^{m^{\text{ans}}}$ for the first 3 equation, and $s \sim \bF_{2^p}^{m}$ for the last equation and over the state $\ket{\tau}$.
\end{claim}
\begin{proof}
$\cG^{\Answerred}$, when the question set is restricted to the case where the oracularization question label is restricted to ``(Prover A)", and the question content for $\cG_n$ is fixed to some $x^{\game} \in \cX$, is precisely an instance of the $(p,m,p)$-low-individual degree test. Since $\strategy$ succeed with probability $1 - c_1 \cdot \delta_1$, and the probability that both oracularization question label in a question pair to be both ``(Prover A)" being $\frac{1}{9}$. This implies that, on average over $(\mu_n)_X$, the strategy $\strategy$, when restricted to the case where $(x^{\text{Ora}}, y^{\text{Ora}})$ being both ``(Prover A)", succeed with probability at least $1 - 9 \cdot \delta_1$. Hence, by~\Cref{thm:soundnessQLDT}
\begin{equation*}
 	A^{(\text{Prover, A}), (\text{Point}), x^{\game} ,s}_{u} \simeq_{\mathbf{\eta}_{\text{LD}}(p,m,p,9 \cdot \delta_1)} G_{[\textbf{eval}_{s}^1 | u]}^{(\text{Prover, A}), x^{\game}},
\end{equation*}
where $\simeq$ for the above equations is defined over the distribution $(x^{\game}) \sim \mu_n$ and $s \sim \bF_{2^p}^{m^{\text{ans}}}$
By the choice of the parameter $p, m^{\text{ans}} = O(\poly(\alpha, \log(n)))$, this immediately shows~\eqref{eq:ARsoundAlicepoly}. ~\eqref{eq:ARsoundBobpoly} and~\eqref{eq:ARsoundvarpoly} follows from a similar argument (except with the oracularization question label replaced with ``(Prover B)" and ``$(\text{Ora})_{o}$", $i = \{0,1,2\}$ respectively). ~\eqref{eq:ARsoundAllpoly} follows a similar argument with the label ``(Ora)" and using~\Cref{lem:simuQLID}. This completes the proof.
\end{proof}
We wish to modify the output for the PVM associated with the ``(Prover, A)", ``(Prover, B)" and ``$(\text{Ora})_{o}$", $i \in \{0,1,2\}$ as PVM which outputs $\textbf{g} \in \Idpoly(p,m,p)$. For $s \in \bF_{2^p}^m$, partition $s = (s_0, \cdots, s_4, w)$ where $s_i \in \bF_{2^p}^{m^{\text{ans}}}$, $i \in [5] $ and $w \in  \bF_{2^p}^{5+g}$, we make the following post measurement processing to PVMs given in the last lemma as follows
\begin{itemize}
	\item Treat the outputs $\textbf{g}$ from $\{G_{\overline{\textbf{g}}}^{(\text{Prover, A}), x^{\game}}\}$ as $\textbf{g} \in \Idpoly(p,m,p)$ as $\textbf{g}(s) =\overline{\textbf{g}}(s_0)$.
	\item Treat the outputs $\textbf{g}$ from $\{G_{\overline{\textbf{g}}}^{(\text{Prover, B}), y^{\game}}\}$ as $\textbf{g} \in \Idpoly(p,m,p)$ as $\textbf{g}(s) =\overline{\textbf{g}}(s_1)$.
	\item For $o \in \{0,1,2\}$, treat the outputs $\textbf{g}$ from $\{G_{\overline{\textbf{g}}}^{(\text{Ora})_{o}, (x^{\game}, y^{\game})}\}$ as $\textbf{g} \in \Idpoly(p,m,p)$ with $\textbf{g}(s) =\overline{\textbf{g}}(s_{o+2})$.
\end{itemize}
To distinguish the two measurement outputs, we write the output polynomial as $\overline{\textbf{g}}$ if the resulting polynomial output are from $\Idpoly(p,m^{\text{ans}},p)$ and $\textbf{g}$ if the output are from $\Idpoly(p,m,p)$. For $i \in [5]$, define the PVM measurement
\begin{align}
	\label{eq:ARproofOracFull-1} G_{\textbf{g}_{U_i}}^{(\text{Ora}), (x^{\game}, y^{\game}), U_i} &= \sum_{\substack{\textbf{g}_{U_0}, \cdots,\textbf{g}_{U_{i-1}},\textbf{g}_{U_{i+1}}, \cdots,\textbf{g}_{U_4} \\ \textbf{g}_{\Gamma}, \textbf{g}_{B_0},\cdots,  \textbf{g}_{B_{m-1}} \in \Idpoly(p, m, p) } }  G_{\textbf{g}_{U_0}, \cdots,\textbf{g}_{U_4},  \textbf{g}_{\Gamma}, \textbf{g}_{B_0},\cdots,  \textbf{g}_{B_{m-1}}  }^{(\text{Ora}), (x^{\game}, y^{\game})} \\
	\label{eq:ARproofOracFull-2} G_{\textbf{g}_{\Gamma}, \textbf{g}_{B_0} \cdots,\textbf{g}_{B_{m-1}}}^{(\text{Ora}), (x^{\game}, y^{\game}), \text{Full}} &= \sum_{\textbf{g}_{U_0}, \cdots, \textbf{g}_{U_4}  \in \Idpoly(p, m, p) }  G_{\textbf{g}_{U_0}, \cdots,\textbf{g}_{U_4},  \textbf{g}_{\Gamma}, \textbf{g}_{B_0},\cdots,  \textbf{g}_{B_{m-1}}  }^{(\text{Ora}), (x^{\game}, y^{\game})}.
\end{align}
Since $ G^{(\text{Ora}), (x^{\game}, y^{\game})} $ is projective, for all $(x,y) \in \cX_n^2$ and outputs $\textbf{g}_v$, $v \in \{U_0, \cdots, U_4, \Gamma, B_0, \cdots, B_{m-1}\}$
\begin{align} \label{eq:ARproofGOrafull}
	 G^{(\text{Ora}), (x, y} = &G^{(\text{Ora}), (x, y),U_0 } \cdots  G^{(\text{Ora}), (x, y), U_4} G^{(\text{Ora}), (x, y), \text{Full}} G^{(\text{Ora}), (x, y),U_4 } \cdots G^{(\text{Ora}), (x, y), U_0}.
\end{align}
Base on the decision procedure for $\cG^{\Answerred}$, we show the following claim 
\begin{claim} \label{claim:ARconsistpoly}
	On average over $(x,y) \sim \mu$, $s \in \bF_{2^p}^m$ and over the state $\ket{\tau}$
	\begin{align}
		\label{eq:ARconsistAlicepoly} G_{[\textbf{eval}_{s}^1 | u_0]}^{(\text{Prover, A}), x^{\game}}	&\simeq_{O(\poly(\eps, \log(n), \alpha))} 	(G_{[\textbf{eval}_{s}^1 | u_0]}^{(\text{Ora}), (x^{\game}, y^{\game}), U_0})^{op}  \\
		\label{eq:ARconsistBobpoly} G_{[\textbf{eval}_{s}^1 | u_1]}^{(\text{Prover, B}), y^{\game}}	&\simeq_{O(\poly(\eps, \log(n), \alpha))} 	(G_{[\textbf{eval}_{s}^1 | u_1]}^{(\text{Ora}), (x^{\game}, y^{\game}), U_1})^{op}  \\
		\label{eq:ARconsistOracpoly} G_{[\textbf{eval}_{s}^{o+2} | u_{o+2}]}^{(\text{Ora})_{o}, (x^{\game}, y^{\game})}	&\simeq_{O(\poly(\eps, \log(n), \alpha))} 	(	G_{[\textbf{eval}_{s}^{o+2} | u_{o+2}]}^{(\text{Ora}), (x^{\game}, y^{\game}), U_{o+2}})^{op} , \, o \in \{0,1,2\}.
	\end{align} 
\end{claim}

\begin{proof}
	We show the proof for~\eqref{eq:ARconsistAlicepoly} below, the proof for~\eqref{eq:ARconsistBobpoly} and~\eqref{eq:ARconsistOracpoly} follows a similar proof. Consider the question pair $(x,y)$ for $\cG^{\Answerred}$ where $(x^{\text{Ora}}, y^{\text{Ora}}) =$ ((Prover A), (Ora)) and $x^{\text{LDL}} = y^{\text{LDL}} = $ (Point), this occur with constant probability. Since $\strategy$ succeed with probability at least $1-  c_1 \cdot\delta_1$, by point 2 of the ``Prover consistency check" from~\Cref{fig:Answerredver},
	\begin{equation} \label{eq:ARproofB1}
		A^{((\text{Prover, A}), (\text{Point}), (x^{\game} ,s))}_{u_0} \simeq_{O(\poly(\eps, \log(n), \alpha))} A^{(\text{Ora}), (\text{Point}), ((x^{\game}, y^{\game}) ,s)}_{u_0, \cdots, u_4, \gamma, \beta_0, \cdots,\beta_{m-1}} 
	\end{equation}
	over $(x,y) \sim \mu$ and $s \sim \bF_{2^p}^m$ and over the state $\ket{\tau}$. ~\eqref{eq:ARconsistAlicepoly} then follows from~\Cref{lem:closetodistance} point 1 and 2 which translates between $\simeq$ distance to $\approx$ distance, and the triangle inequality for $\approx$ distance applied to~\eqref{eq:ARsoundAlicepoly},~\eqref{eq:ARproofB1}, and~\eqref{eq:ARsoundAllpoly}. 
\end{proof}

For $(x^{\game}, y^{\game}) \in \cX_n$, define the POVM measurement $M^{(x^{\game}, y^{\game})}_{\textbf{g}_{U_0}, \cdots, \textbf{g}_{U_4}, \textbf{g}_{\Gamma}, \textbf{g}_{B_0} \cdots,\textbf{g}_{B_{m-1}}}$ with outcomes $\textbf{g}_v \in \Idpoly(p,m,p)$ as
\begin{align*}
	M^{(x, y)}_{\textbf{g}_{U_0}, \cdots, \textbf{g}_{U_4}, \textbf{g}_{\Gamma}, \textbf{g}_{B_0} \cdots,\textbf{g}_{B_{m-1}}} = &G_{\textbf{g}_{U_0}}^{(\text{A}), x} G_{\textbf{g}_{U_1}}^{(\text{B}), y}  G_{\textbf{g}_{U_2}}^{(\text{Orc})_0, (x,y)}  G_{\textbf{g}_{U_3}}^{(\text{Orc})_1, (x,y)}  G_{\textbf{g}_{U_4}}^{(\text{Orc})_2, (x,y)} 	G_{\textbf{g}_{\Gamma}, \textbf{g}_{B_0} \cdots,\textbf{g}_{B_{m-1}}}^{(\text{Ora}), (x, y), \text{Full}} \cdot \\
	&\quad G_{\textbf{g}_{U_4}}^{(\text{Orc})_2, (x,y)}  G_{\textbf{g}_{U_3}}^{(\text{Orc})_1, (x,y)}  G_{\textbf{g}_{U_2}}^{(\text{Orc})_0, (x,y)} G_{\textbf{g}_{U_1}}^{(\text{B}), y} G_{\textbf{g}_{U_0}}^{(\text{A}), x}, 
\end{align*}
where for $P \in \{A,B\}$, we shorten the label (Prover, P) to (P), and remove the superscript ``\game" in the above equation. We remark that in contrast to $G^{(\text{Ora}), (x^{\game}, y^{\game})}$, the output $\textbf{g}_{U_i}$, $i \in [5]$ from $M^{(x, y)}$ are secretly polynomials in $\Idpoly(p, m^{\text{ans}},m)$ which is consistent with the definition for the 5 polynomials given in~\Cref{thm:classicalPCPP}. 

Furthermore, we define $M^{(x^{\text{game}}, y^{\text{game}}), U_i}$ for $i \in [5]$, and $M^{(x^{\text{game}}, y^{\text{game}}), \text{Full}}$ in a similar manner as~\eqref{eq:ARproofOracFull-1} and~\eqref{eq:ARproofOracFull-2}. By definition
\begin{align*}
	M^{(x^{\text{game}}, y^{\text{game}}), U_0} &=  G^{(\text{Prover, A}), x^{\game}} \\
	M^{(x^{\text{game}}, y^{\text{game}}), U_1} &=  G^{(\text{Prover, B}), y^{\game}} \\
	M^{(x^{\text{game}}, y^{\text{game}}), U_{o+2}} &=  G_{\overline{\textbf{g}}}^{(\text{Ora})_{o}, (x^{\game}, y^{\game})}, o \in \{0,1,2\}.
\end{align*}
For $(x^{\text{game}}, y^{\text{game}}) \in \cX^2$ and output tuple $(\textbf{g}_{U_0}, \cdots, \textbf{g}_{U_4}, \textbf{g}_{\Gamma}, \textbf{g}_{B_0} \cdots,\textbf{g}_{B_{m-1}})$ from $M^{(x^{\text{game}}, y^{\text{game}})}$. We refer to the output as ``good" if for a uniformly random $s =(s_0, \cdots, s_4, b_0, \cdots, b_4, z) \sim \bF_{2^p}^m$, the following occurs with probability over $\frac{1}{2}$:
\begin{itemize}
	\item $\textbf{g}_{\Gamma}(s) = \textbf{g}_{\deciderTM}(s) (\textbf{g}_{U_0}(s) - b_0)  (\textbf{g}_{U_1}(s) - b_1)  (\textbf{g}_{U_2}(s) - b_2)  (\textbf{g}_{U_3}(s) - b_3)  (\textbf{g}_{U_4}(s) - b_4)$
	\item $\textbf{g}_{\Gamma}(s) = \sum_{i \in [m]} \textbf{g}_{B_i}(s) \textbf{zero}(s) $,
\end{itemize}
where $\textbf{g}_{\deciderTM} = \mathtt{ComputePCP}_{\alpha}(\langle \deciderTM \rangle ,n, x^{\text{game}}, y^{\text{game}})$. By~\Cref{thm:classicalsoundnessPCPP}, if the output for $M^{(x^{\text{game}}, y^{\text{game}})}$ is a ``good" output, then there exist $a,b \in \{0,1\}^*$ with $|a|, |b|\leq \log^{\alpha}(n)$ such that $\textbf{g}_{U_0} =  \text{enc}_{\Gamma}(a)$ and $\textbf{g}_{U_1} =  \text{enc}_{\Gamma}(b)$ and $D(x^{\text{game}}, y^{\text{game}},a,b) =1$

We now proof the following claim regarding the measurement $M^{(x^{\text{game}}, y^{\text{game}})}$
\begin{claim} \label{claim:ARgoodoutput}
	On average over $(x^{\text{game}}, y^{\text{game}}) \sim \mu_n$ and the state $\ket{\tau}$
	\begin{equation} \label{eq:ARproofSynccondpoly}
		M^{(x^{\text{game}}, y^{\text{game}})} \simeq_{O(\poly(\eps, \log(n), \alpha))}  (M^{(x^{\text{game}}, y^{\text{game}})})^{op}.
	\end{equation}
	Furthermore, on average over $(x^{\text{game}}, y^{\text{game}})$, the measurement output for $\braket{\tau|M^{(x^{\text{game}}, y^{\text{game}})}|\tau}$ is ``good" with probability at least $1 - O(\poly(\eps, \log(n), \alpha))$
\end{claim}
\begin{proof}
	 We first show that, on average over $(x^{\text{game}}, y^{\text{game}}) \sim \mu_n$ and the state $\ket{\tau}$
	\begin{equation} \label{eq:ARsoundnessgoodoutcome-1}
		M^{(x^{\text{game}}, y^{\text{game}})} \simeq_{O(\poly(\eps, \log(n), \alpha))}  G^{(\text{Ora}), (x^{\game}, y^{\game})}
	\end{equation}
	Since $G^{(\text{Ora}), (x^{\game}, y^{\game})}$ is projective, by the definition of $M^{(x^{\text{game}}, y^{\text{game}})}$, the last $1+m$ measurement outcome for $G^{(\text{Ora}), (x^{\game}, y^{\game})}$ ($ \textbf{g}_{\Gamma}, \textbf{g}_{B_0},\cdots,  \textbf{g}_{B_{m-1}}$) will always be the same as the measurement $M^{(x^{\text{game}}, y^{\text{game}})}$ when the two measurements are made simultaneously (on any state). 
	For $i \in [5]$, by~\Cref{claim:ARconsistpoly} and the Schwartz-Zippel lemma (\Cref{lem:Schwartz_Zipple})
	\begin{align*}
		M^{(x^{\text{game}}, y^{\text{game}}), U_i}_{\textbf{g}_{U_i}}&\simeq_{\delta_2} (G_{\textbf{g}_{U_i}}^{(\text{Ora}), (x^{\game}, y^{\game}), U_i})^{op}
	\end{align*}
	for $\delta_2 = O(\poly(\eps, \log(n), \alpha)) + \frac{m \cdot d}{2^p}$. Hence, by repeatedly applying~\Cref{lem:combmeasure}, the underlying vector state is a tracial state, and using~\Cref{lem:closetodistance} to convert between $\simeq$ distance to $\approx$ distance,
	\allowdisplaybreaks{
		\begin{align*}
			 &G^{(\text{Ora}), (x, y)} = G^{(x, y),U_0 } \cdots  G^{(x, y), U_4} G^{(x, y), \text{Full}} G^{(x, y),U_4 } \cdots G^{ (x, y), U_0} \\
			 &\approx_{\delta_2} (M^{(x, y), U_0})^{op} G^{(x, y),U_0 } \cdots  G^{(x, y), U_4} G^{(x, y), \text{Full}} G^{(x, y),U_4 } \cdots G^{(x, y),U_1 } \\
			 &\cdots \\
			 &\approx_{\delta_2} (M^{(x, y), U_4} \cdots M^{(x, y), U_0})^{op}  G^{ (x, y),U_0 } \cdots  G^{(x, y),U_4 }   G^{(x, y), U_4} G^{(x, y), \text{Full}}  \\
			 &=  (M^{(x, y), \text{Full}} M^{(x, y), U_4} \cdots M^{(x, y), U_0})^{op}  G^{ (x, y),U_0 } \cdots  G^{(x, y),U_4 } \\
			 &\approx_{\delta_2}  (M^{(x, y), U_4} M^{(x, y), \text{Full}} M^{(x, y), U_4} \cdots M^{(x, y), U_0})^{op}  G^{ (x, y),U_0 } \cdots  G^{(x, y),U_3 }  \\
			 &\cdots \\
			 &\approx_{\delta_2}  ( M^{(x, y), U_0} \cdots M^{(x, y), U_4} M^{(x, y), \text{Full}} M^{(x, y), U_4} \cdots M^{(x, y), U_0})^{op}  =  (M^{(x, y)})^{op}
		\end{align*}
	}
	where we remove the superscript ``game" and (\text{Ora}) in the above derivation for clarity. Hence, by using the triangle inequality for $\approx$ distance and~\Cref{lem:closetodistance}, this implies that 
	\begin{equation*}
		G^{(\text{Ora}), (x^{\text{game}}, y^{\text{game}})}  \simeq_{10 \delta_2}(M^{(x^{\text{game}}, y^{\text{game}})})^{op}
	\end{equation*}
	and since the underlying state is a tracial state and $\delta_2 = O(\poly(\eps, \log(n), \alpha))$, this shows~\eqref{eq:ARsoundnessgoodoutcome-1}. Since the underlying state for~\eqref{eq:ARsoundnessgoodoutcome-1} is the tracial state $\ket{\tau}$, we also have
	\begin{equation} \label{eq:ARsoundnessgoodoutcome-2}
		(M^{(x^{\text{game}}, y^{\text{game}})})^{op} \simeq_{O(\poly(\eps, \log(n), \alpha))}  (G^{(\text{Ora}), (x^{\game}, y^{\game})})^{op}
	\end{equation}
	For~\eqref{eq:ARproofSynccondpoly}, since $G^{(\text{Ora}), (x^{\game}, y^{\game})} $ are all projective measurements, 
	\begin{equation} \label{eq:ARsoundnessgoodoutcome-3}
		G^{(\text{Ora}), (x^{\game}, y^{\game})} \simeq_0 (	G^{(\text{Ora}), (x^{\game}, y^{\game})})^{op}
	\end{equation}
	over $(x^{\game}, y^{\game}) \sim \mu_n$ and the tracial state $\ket{\tau}$.~\Cref{eq:ARproofSynccondpoly} then follows from~\Cref{lem:closetodistance} and the triangle inequality of $\approx$ distance being applied to \eqref{eq:ARsoundnessgoodoutcome-1}, \eqref{eq:ARsoundnessgoodoutcome-2} and \eqref{eq:ARsoundnessgoodoutcome-3}. 
	
	For the second part of~\Cref{claim:ARgoodoutput}, by applying the data processing inequality to~\eqref{eq:ARsoundnessgoodoutcome-1},
	\begin{equation} \label{eq:ARsoundnessgoodoutcome-4}
		M^{(x^{\text{game}}, y^{\text{game}})}_{[\textbf{eval}_s^{6+m}|( u_0, \cdots, u_4, \gamma, \beta_0, \cdots,\beta_{m-1})]} \simeq_{O(\poly(\eps, \log(n), \alpha))} A^{(\text{Ora}), (\text{Point}), ((x^{\game}, y^{\game}) ,s)}_{[\textbf{eval}_s^{6+m}|( u_0, \cdots, u_4, \gamma, \beta_0, \cdots,\beta_{m-1})]}.
	\end{equation}
	Hence by applying the triangle inequality for $\approx$ distance and~\Cref{lem:closetodistance} to~\eqref{eq:ARsoundnessgoodoutcome-4} and~\eqref{eq:ARsoundAllpoly}, we obtain
	\begin{equation}~\label{eq:ARsoundnessgoodoutcome-5}
		M^{(x^{\text{game}}, y^{\text{game}})}_{[\textbf{eval}_s^{6+m}|( u_0, \cdots, u_4, \gamma, \beta_0, \cdots,\beta_{m-1})]} \simeq_{O(\poly(\eps, \log(n), \alpha))}  A^{(\text{Ora}), (\text{Point}), ((x^{\game}, y^{\game}) ,s)}_{u_0, \cdots, u_4, \gamma, \beta_0, \cdots,\beta_{m-1}}. 
	\end{equation}
	Recall that $\strategy$ is a synchronous strategy which succeed at $\cG^{\text{AR}}$ with probability at least $1 -  c_1 \cdot \delta_1$. Since the oracularization question label is pick with constant probability, by the ``PCPP proof check" clause of~\Cref{fig:Answerredver}, on expectation over $(x^{\game},y^{\game}) \sim \mu$ and $s = (s_0, \cdots, s_4, b_0, \cdots, b_4, z) \in \bF_{2^p}^m$, the measurement 
	\begin{equation*}
		\braket{\tau|A_{(u_0, \cdots, u_4, \gamma, \beta_0, \cdots, \beta_{m-1} )}^{(\text{Orc}), (\text{Point}), ((x^{\game}, y^{\game}),  s)}|\tau} 
	\end{equation*}	
	outputs the answer which satisfies the properties below with probability $1 - c_2 c_1 \cdot\delta_1$
	\begin{enumerate}
		\item $\gamma = \textbf{g}_{\deciderTM}(s) \cdot (u_1 - b_0) \cdots (u_4 - b_4)$,
		\item $\gamma = \sum_{i \in[m]} \beta_i \cdot \textbf{zero}(s_i)$,
	\end{enumerate}
	where $\textbf{g}_{\deciderTM} = \mathtt{ComputePCP}_{\alpha}-(\langle \deciderTM \rangle ,n, x^{\text{game}}, y^{\text{game}})$. Pick $c_1 \in (0,1)$ used to define $\strategy$ such that $1 - c_2 c_1 \cdot\delta_1 \geq \frac{1}{2}$. Combine this with~\Cref{eq:ARsoundnessgoodoutcome-5}, this shows that on average over $(x^{\game}, y^{\game}) \sim \mu_n$, the probability that $M^{(x^{\text{game}}, y^{\text{game}})}$ gives an output which is ``good" with probability at least $1-= O(\poly(\eps, \log(n), \alpha))$, thus completing the claim for the lemma. 
\end{proof}
Base on the POVM $M^{(x^{\text{game}}, y^{\text{game}})}$, we define a symmetric strategy $\strategy^{\text{Ora}} =(\cL^2(\alicealg, \tau), \ket{\tau}, \{B_a^x \})$ for $\cG^{\text{Ora}}$ as follows: Fixed $(x^{\text{game}}, y^{\text{game}}) \in \cX_n$, the measurement operator $\{B^{\text{(Prover, A)}, x^{\text{game}}}_a\}$ as a data processing measurement as follows:
\begin{itemize}
	\item Perform the measurement $M^{x^{\text{game}}, U_0}_{\textbf{g}_{U_0}}$ and obtain a polynomial $\textbf{g}_{U_0}$.
	\item If there exist an $a \in \cA_n$ such that $|a| \leq \log^{\alpha}(n)$ and $\textbf{g}_{U_0} = \text{enc}$, output $a$. Otherwise, output $0$. 
\end{itemize}
The measurement operator $\{B^{\text{(Prover, B)}, y^{\text{game}}}\}$ is define in a similar manner as $\{B^{\text{(Prover, A)}, x^{\text{game}}}\}$ with the measurement $M^{x^{\text{game}}, U_1}_{\textbf{g}_{U_1}}$. The measurement operator $\{B^{\text{(Orac)}, (x^{\text{game}},y^{\text{game}})}\}$ similarly define as a data processing measurement as follows:
\begin{itemize}
	\item Perform the measurement $M^{(x^{\text{game}}, y^{\text{game}})}$ and obtain the tuple of polynomials $(\textbf{g}_{v})$.
	\item If the given measurement outcome is a ``good" output, by~\Cref{thm:classicalsoundnessPCPP}, there exist $(a,b) \in \cA_n^2$ such that $\textbf{g}_{U_0} = \text{enc}_{\Gamma}(a)$, $\textbf{g}_{U_1} = \text{enc}_{\Gamma}(b)$, output the corresponding $(a,b) \in \cA_n^2$. Otherwise, output $(0,0)$. 
\end{itemize}
We remark that the above strategy is not necessarily synchronous strategy, since $\{M^{(x^{\text{game}}, y^{\text{game}})}\}$ does not necessarily have to be a PVM. We make the following claim about $\strategy^{\text{Ora}}$. 
\begin{claim}
	$\omega(\cG^{\textbf{Orac}}, \strategy^{\text{Ora}}) \geq 1 - O(\poly(\eps, \log(n), \alpha))$. 
\end{claim}
\begin{proof}
	We consider the performance of $\strategy^{\text{Ora}}$ for different question pairs given in~\Cref{fig:Oracularization} below:
	\begin{itemize}
		\item (Prover, P) \textbf{-}  (Prover, P), $P \in \{A,B\}$: Since both $M^{x^{\text{game}}, U_A} = G^{(\text{Prover, A}), x^{\game}} $ and $M^{y^{\text{game}}, U_B} = G^{(\text{Prover, B}), y^{\game}} $ are both projective and $\strategy^{\text{Ora}}$ uses the tracial state as the underlying state. This implies that $\strategy^{\text{Ora}}$ always succeed on this question pair.
		\item 	(Oracularization) \textbf{--} (Oracularization): For the ``consistency" part of this question pair,  $B^{\text{(Orac)}, (x^{\text{game}},y^{\text{game}})}$ is a data processed measurement of $M^{(x^{\text{game}}, y^{\text{game}})}$, which by~\eqref{eq:ARproofSynccondpoly}, are consistent with probability at least $1 - O(\poly(\eps, \log(n), \alpha))$. For the ``proof checking" part of this question pair,, whenever $M^{(x^{\text{game}}, y^{\text{game}})}$ returns a ``good" output when performing the measurement $B^{\text{(Orac)}, (x^{\text{game}},y^{\text{game}})}$, the corresponding output $(a,b) \in \cA_n^2$ always satisfies $D_n(x,y,a,b) = 1$ by~\Cref{thm:classicalsoundnessPCPP}. By~\Cref{claim:ARgoodoutput}, this occurs with $1 - O(\poly(\eps, \log(n), \alpha))$. Combining these two facts, this implies that $\strategy^{\text{Ora}}$ succeed on this question pair with probability at least $1 - O(\poly(\eps, \log(n), \alpha))$.  
		\item (Oracularization) \textbf{--} (Prover, P), $P \in \{A,B\}$: Restricted to the case when the provers receiving the question label `` (Oracularization)" and obtain a ``good" outcome from the measurement of $M^{(x^{\text{game}}, y^{\text{game}})}$, by the definition of $M$, the `` (Prover, P)" prover would receive the same polynomial from his/her measurement output, and hence output a consistent answer label as the `` (Oracularization)" prover. Since a ``good" outcome occurs with probability  $1 - O(\poly(\eps, \log(n), \alpha))$, this implies that $\strategy^{\text{Ora}}$ succeed with probability on this question pair with probability at least $1 - O(\poly(\eps, \log(n), \alpha))$.  
	\end{itemize}
	By averaging out the probability given above, we see that $\omega(\cG^{\textbf{Orac}}, \strategy^{\text{Ora}}) > 1 - O(\poly(\eps, \log(n), \alpha))$, completing the proof of the claim. 
\end{proof}
This shows that for model $t \in \{*,co\}$, $\omega^t(\cG) > 1- \eps$ implies that $\omega^t(\cG^{\text{Ora}}) > 1- O(\poly(\eps, \log(n), \alpha))$. The proof of~\Cref{prop:AnswerReduction} then follows from the ``soundness" clause of~\Cref{lem:oratransformation}. 


% Fix $(x^{\game}, y^{\game}) \in \cX_n$, for $P \in \{A,B\}$ and $x^{\text{game}} \in \cX_n$, we define the measurement operator $\{B^{\text{(Prover, P)}, x^{\text{game}}}_a\}$ as follows: 



\begin{comment}
Since $\strategy$ is synchronous, projective strategy. By~\Cref{lem:closetodistance}, we can rewrite~\Cref{eq:ARsoundAlicepoly} as 
\begin{equation*}
	\left(A^{(\text{Prover, A}), (\text{Point}), x^{\game} ,s}_{u}\right)^{\text{op}} \approx_{O(\poly(\eps, \log(n), \alpha))} \left(G_{[\textbf{eval}_{s}^1 | u]}^{(\text{Prover, A}), x^{\game}} \right)^{\text{op}} , 
\end{equation*}
and since $A^{(\text{Prover, A}), (\text{Point}), x^{\game} ,s}_{u} \approx_{0} \left(A^{(\text{Prover, A}), (\text{Point}), x^{\game} ,s}_{u}\right)^{\text{op}}$ (because $\strategy$ is synchronous), by the triangle inequality of $\approx$ and as well as point 2 of~\Cref{lem:closetodistance}, 
\begin{equation}
	G_{[\textbf{eval}_{s}^1 | u]}^{(\text{Prover, A}), x^{\game}} \approx_{O(\poly(\eps, \log(n), \alpha))} \left(G_{[\textbf{eval}_{s}^1 | u]}^{(\text{Prover, A}), x^{\game}} \right)^{\text{op}},
\end{equation}
and since 


and similar statement can be made about all the PVMs guaranteed by~\Cref{claim:ARpolyPVM}. 
\end{comment}


\begin{comment}
	The third lemma allows us to combine data processing measurements, we remark that this is an analogue of~\cite[Fact 4.34]{natarajanNEEXPMIP2019a}. 
\begin{lemma}[Combination of low-individual degree polynomials]
	Fix constant $p, m, d, k \in \bN$, finite set $\cX$ and some distribution $\mu$ over $\cX$. For each $i \in [k]$ and $x \in \cX$, let $\{G_{\textbf{g}}^{x,i}\}$ be a set of POVM with outcomes $\textbf{g} \in \Idpoly(p,m,d)$ being the sets of low-individual degree polynomials of individual degrees of at most $d$. Let $\{ H_{\textbf{h}_0, \cdots, \textbf{h}_{k-1}}^x\}$ be a set of POVM with outcomes being $\textbf{h}_i \in \Idpoly(p,m,d)$ for each $i \in [k]$. For each $i \in [k]$, write
	\begin{equation*}
		H_{\textbf{h}_{i}}^x = \sum_{\textbf{h}_{0} \cdots, \textbf{h}_{i-1}, \textbf{h}_{i+1}, \cdots \textbf{h}_{k-1} }  H_{\textbf{h}_0, \cdots, \textbf{h}_{k-1}}^x
	\end{equation*}
	and suppose that 
	\begin{equation}
		H_{[\textbf{h}_i| a]} \simeq_{\delta} 
	\end{equation}
	
\end{lemma}

\end{comment}

